\honest{Zombies! Zombies?}{Eliezer Yudkowsky}{2008}

Your zombie, in the philosophical sense, is exactly like you down to the last atom, moving under the same laws, except that it is not conscious. It is claimed that if such zombies are possible, then purely from that possibility we can deduce that consciousness is beyond the physical, and the standard term for this position is ``epiphenomenalism'': consciousness exists but does nothing. \nb{The zombie argument concludes only that physicalism is false. Epiphenomenalism needs the further premise that the physical world is causally closed, and the \textsc{Stanford Encyclopedia of Philosophy} notes that friends of zombies may instead be interactionists.} This is not a strawman: ``possibly a majority'' of philosophers of consciousness accept it. \nb{In the 2009 survey, a year later, 18.3 per cent of philosophers of mind held zombies metaphysically possible and 47.6 per cent held them conceivable but not possible.} This long post engages a real philosopher, David Chalmers, in answer to the charge that I do not.

I once read that ``You are not the one who speaks your thoughts'': ``you are the one who hears your thoughts.'' In Hebrew the highest soul, the one God breathed into Adam, is the N'Shama, ``the hearer.'' If consciousness is a passive listener, a zombie is easy to imagine: someone without the hearer.

When you find the orange juice gone and think ``Darn, I'm out of orange juice,'' the sounds are probably represented in your auditory cortex. (Chinese speakers remember about ten digits, against the famous seven plus or minus two for English speakers, because Chinese digits are single syllables; working memory seems to be a loop of sounds.) Within decades surgeons may read this inner narrative with neural taps; researchers have already reconstructed images from a cat's lateral geniculate nucleus. Your zombie forms the same patterns for the same words, and epiphenomenalists agree. But, they say, in the zombie no one hears them. And since even knowing the position of every atom you would still have to be told, as a further fact, that people have inner listeners, consciousness must be something over and above atoms.

This is not dualism. Descartes's mind-substance was causally active and controlled speech and behaviour; subtract it and you get a traditional zombie, lurching and groaning. And I cannot recall a rabbi arguing for zombies; most would be aghast that the part God breathed into Adam does nothing. The passive listener, ``the lights are on but no one's home,'' is what first seduces people, especially lay audiences, to the zombie view.

The argument needs the zombie world to be logically possible, not just to look possible. In logic a set of statements such as $A \to B$, $B \to C$, $C \to \neg A$ is possible if it has a model, here an assignment of true or false that makes all three true ($A=B=C=$ false works). Something seems possible if you consider it without seeing a contradiction. But finding contradictions, or models, is hard in general: for clauses of three variables the problem is 3-SAT, one of the first problems proven NP-complete. Not seeing a contradiction in the zombie world at first glance is like not seeing one in the Riemann Hypothesis. You might see it in another thirty seconds: ``All odd numbers are prime. Proof: 3 is prime, 5 is prime, 7 is prime\ldots'' \nb{Chalmers makes the same distinction in \textsc{The Conscious Mind} (1996), and gives arguments that zombies pass the stricter test. The post does not take those up.}

So ponder a little longer. Is there a counterexample to ``Consciousness has no third-party-detectable causal impact on the world''? Close your eyes and attend to your awareness, and your inner voice says ``I am aware,'' ``I am not the one who speaks my thoughts, but the one who hears them,'' and so on. Say it aloud and your lips move. The listener is ``caught in the act of listening.'' Compare a black box left by aliens. You poke it and it does nothing, no gold, no answers, so you conclude that for all X the box does not do X. But you can see it, and it is heavy in your hand; that too is causality. If it were wholly outside the causal universe you could not know it existed or thank the aliens for it. (They left a second, truly epiphenomenal box in your living room, and you have no idea it is there. That was their joke.) I have not seen this argument put this way, ``though it may well have been said before.'' \nb{It had been: by G.~F. Stout in 1931, and by Avshalom Elitzur in 1989, whom Chalmers credits in the book this post quotes.}

Zombie advocates accept that zombie philosophers, being atom for atom identical to ours, write the same papers about consciousness. At that point the zombie world stops following from the passive listener. Papers about consciousness look like an effect of consciousness; it is usually right after you attend to your awareness that you say you are aware. You can argue cleverly that this is not so, but you have to be clever.

Could a Zombie Master, a god in the zombie world, make the zombies talk about consciousness? Faking such talk is not hard; human talk about consciousness is not very coherent, and an AI trained on a thousand amateurs in a bar might match it. But that would be imitation, a holodeck, not what the zombie world is. The zombie world has the same atoms under the same laws, and any controller of lips would be detectable: some electron would zag instead of zig because the Master said so. When a philosopher types Z-O-M-B-I-E, the causal chain runs back through muscles, nerves, spinal cord and motor cortex to the part of the brain where the inner narrative began to talk about consciousness, and the zombie twin types the same keys for the same physical reasons. So you cannot say the philosopher writes because of consciousness and the twin because of a Zombie Master.

I quote Chalmers on his zombie twin, who is obsessed with consciousness, writes chapter after chapter on it, professes a love of deep greens and purples, argues with zombie materialists, and is wrong; yet there must be a physical explanation of his claims, and any explanation of the twin's behaviour also explains Chalmers's, so the explanation of his own claims ``is also independent of the existence of consciousness.'' These are Chalmers's actual beliefs. Chalmers finds this ``at once delightful and disturbing''; I call it ``the largest bullet ever bitten in the history of time,'' and that is a backhanded compliment: a lesser mortal would not have seen the implication.

Why bite it? Not because of the passive listener, which may let zombies drive or fall in love but not write papers about their listeners. The drive comes from a second intuition: no sum of masses and charges will necessarily produce the redness of red, so if certain atoms evoke redness, that is an extra, contingent fact. But that intuition alone does not give zombies. Descartes's substance dualism also has a different kind of stuff, but one that acts: remove it and you stop writing papers about consciousness. Zombie advocates are property dualists: matter has extra properties that do not move atoms, so no third party can detect them, which is why we can supposedly imagine them gone with everything else the same. The zombie world, they say, is not physically possible, but it is logically possible, since the bridging laws could have been different.

Then why say the extra properties do nothing? Then you need both an extra-physical redness of red and a wholly separate physical reason why Chalmers talks about the redness of red. Descartes's view is strictly simpler: once you posit a mysterious stuff of consciousness, let it make you talk about consciousness. \nb{\textsc{The Conscious Mind} says: ``I do not describe my view as epiphenomenalism,'' and leaves the causal question open, while granting ``at least a weak form.'' By 2003 Chalmers gave credence to interactionism and panprotopsychism as well. He also says the paradox arises ``almost as strongly'' on the panprotopsychist view, and that interactionism is ``of little help'' against it, so the argument survives the correction.} Vitalists would spit out their coffee at the alternative: ``When fires burn, they release phlogiston. But phlogiston doesn't have any experimentally detectable impact on our universe, so you'll have to go looking for a separate explanation of why a fire can melt snow.'' If you posit an extra mental property, you have already stuck your neck out as far as it goes. There is not even a career motive for saying the thing you study does nothing. Chalmers objects that it is hard to see what new physics could explain consciousness; property dualism has the same problem. \nb{Chalmers gives other reasons: current evidence that physics is causally closed, and, as his ``deepest reason,'' that the phenomenal can be subtracted from any causal story, even one with immaterial ``psychons.'' He agrees the problem is the same, and concludes that interaction gains nothing.}

An unkind critic would bring in Carl Sagan's invisible, inaudible dragon in the garage, which does not breathe and is permeable to flour: one motive for making a theory unfalsifiable is fear of testing it. Penrose and Hameroff are substance dualists who think consciousness lives in a physically real collapse of the wave-function, and they predicted that neurons keep macroscopic quantum states; tests so far look bad for them, but their conduct is scientifically respectable. \nb{In a 2014 review they place their theory under ``Science,'' not ``Dualism.''} I once told Hameroff that his hypothesis was completely hopeless and his experiments should definitely be funded. But I do not think fear of testing is true of Chalmers. (And if he had it, he should note that if epiphenomenalism is false, the real explanation will eventually be found and ruin his theory.)

Chalmers is one of the most frustrating philosophers I know: he analyses sharply and then turns left at the last minute, reducing the zombie argument to smithereens and then accepting it. He does the same with justification. On his theory his belief in consciousness is not caused by consciousness, since he would write the same papers without it; and it is not the output of a reliable process, since his twin uses the same process and is wrong. Chalmers explains this himself. Then he answers that having the experiences is what justifies his beliefs, and that the question is whether he is justified, not whether his brain is. \nb{Chalmers answers ``But your zombie twin would say the same thing!'' two paragraphs later in the same book: ``So what?''} My reply: Chalmers wrote that in a physical book, and so did his twin, for physical reasons. So even if there is a conscious inner Chalmers, there is a causally closed outer Chalmers that speaks the inner narrative and writes the papers. On Chalmers's own theory the outer Chalmers is deranged: it writes ``for no valid reason,'' and is right only ``by a separate and additional miracle,'' contingent, true in our universe though science could never tell it from the zombie world. \nb{Chalmers holds that the match comes from laws of nature, ``no more miraculous than any contingent law,'' though in 2003 he granted it seems ``something of a lucky coincidence.''} It is like saying that epicycles do nothing, something else moves the planets as the epicycles say, and I would say all this even if there were no epicycles.

I look at philosophy as someone who wants to design a self-improving AI with stable goals, and for that, reflective coherence is essential. If such an AI finds a part of itself that writes ``B'' to memory under condition A, studies how that part works causally, and finds it systematically writes false data, it will treat it as a bug and change it. An AI scanning Chalmers would see a closed causal system producing nonsense, nonsense that matters to what Chalmers thinks counts as a morally valuable person. \nb{Chalmers grants that from the outside he and his zombie twin cannot be told apart. His claim is about justification from the inside, which inspection does not reach.} This is only a problem if epiphenomenalism is true; on reductionism or substance dualism, people who talk about consciousness talk about something real, and an AI can see so by tracing the causes of their words. On Chalmers's view, the narrative malfunctions yet happens to be right in our universe, and even its statement that it malfunctions but happens to be right happens to be right. ``Oh, come on!''

Spelled out: the theory posits several unexplained miracles, which lowers its prior by the conjunction rule and Occam's razor. It is dominated by substance dualism, a stuff of consciousness that visibly affects the world and makes us talk about it, and by reductionism, on which consciousness happens within physics in a way not yet understood, as with the last three thousand mysteries, and your intuition that matter cannot add up to it is mistaken. Epiphenomenalism needs extra properties that move no particles, a separate physical reason for the talk, and a miracle by which our bridging laws make that malfunctioning talk correct. \nb{Chalmers's reply is that the rivals do not cover the same data, since experience is known directly, and he grants that without that evidence the simpler view would win. That premise is the point in dispute.} The zombie argument ``may be a candidate for the most deranged idea in all of philosophy.'' Relativity, quantum mechanics and natural selection seem weird too, but massive evidence pins them down; here there is only a complicated argument full of blank spots, which leaves consciousness as mysterious as before even if you accept it all. A rationalist should say: I don't know how consciousness works, but this ``can't possibly be right.'' This intuition does not replace careful refutation; System 1 can point the way, but you still have to find the specific flaws.

I have not read all of Chalmers's book or retraced his long arguments. I answer only his last defense, which he ``has not yet countered to my knowledge.'' \nb{In 2018 Chalmers quoted this post and replied to arguments of its kind. He conceded that ``it is hard to avoid a sense of coincidence entirely,'' and that most non-epiphenomenalist views face ``a weaker form of the same critique.''}
